\section{Transience and Recurrence}

\begin{definition}[Transience, Recurrence, Return Times, Return Probabilities]
    Suppose \(X\) is a Markov chain on \(I\) with transition matrix \(P\).
    \begin{itemize}
        \item A state \(i \in I\) is \term{recurrent}, if \(\Prob_i \pqty{X_n = i \text{ for infinitely many } n} = 1\).
        \item The transition matrix \(P\) is \term{recurrent}, if all states are recurrent.
        \item A state \(i \in I\) is \term{transient}, if \(\Prob_i \pqty{X_n = i \text{ for infinitely many } n} = 0\).
        \item The transition matrix \(P\) is \term{transient}, if all states are transient.
        \item Define the \term{return times} to \(i\) as
        \[
            T_i^{\pqty{0}} = 0 \qand T_{i}^{\pqty{k + 1}} = \inf \set{n > T_i^{\pqty{k}}}{X_n = i} \qfor k \geq 0
        \]
        and the \term{first return time} as \(T_i^{\pqty{1}} = T_i\).
        \item Define the \term{return probabilities} as \(f_i = \Prob_i \pqty{T_i < \infty}\).
        \item The \term{total number of visits to \(i\)} is given by \(V_i = \sum_{k = 0}^{\infty} \indicator \Bqty{X_k = i}\).
    \end{itemize}
\end{definition}

\begin{lemma}[Return Probabilities and Total Number of Visits]
    \label{lem:return-probabilities-and-total-number-of-visits}%
    For every \(r \geq 0\), \(\Prob_{i} \pqty{V_i > r} = f_i^r\).
\end{lemma}

\begin{proof}
    We prove by induction on \(r\). This is obvious for \(r = 0\), since
    \[
        \Prob_i \pqty{V_i > 0} = 1 = f_i^0.
    \]

    Suppose this is true for \(r\). Then we have
    \begin{align*}
        \Prob_i \pqty{V_i > r + 1} &= \Prob_i \pqty{V_i > r + 1, V_i > v} \\
        &= \Prob_i \pqty{V_i > r} \cdot \Prob_i \pqty{\Conditional{V_i > r + 1}{V_i > r}} \\
        &= f_i^r \Prob_i \pqty{\Conditional{T_i^{\pqty{r + 1}} < \infty}{T_i^{\pqty{r}} < \infty}} \\
        &= f_i^r \Prob_i \pqty{T_i < \infty} \\
        &= f_i^{r + 1}.
    \end{align*}
\end{proof}

\begin{theorem}[States are either Recurrent or Transient]
    Let \(X\) be a Markov chain on \(I\).

    \begin{enumerate}
        \item If \(f_i = \Prob_i \pqty{T_i < \infty} = 1\), then \(i\) is recurrent, and
        \[
            \sum_{n \geq 0} p_{ii} \pqty{k} = \infty.
        \]

        \item If \(f_i = \Prob_i \pqty{T_i < \infty} < 1\), then \(i\) is transient, and
        \[
            \sum_{n \geq 0} p_{ii} \pqty{k} < \infty.
        \]
    \end{enumerate}
\end{theorem}

\begin{proof}
    First notice that
    \begin{align*}
        \Expt_i \pqty{V_i} &= \Expt_i \pqty{\sum_{k = 0}^{\infty} \indicator \Bqty{X_k = i}} \\
        &= \sum_{k = 0}^{\infty} \Prob_i \pqty{X_k = i} \\
        &= \sum_{k = 0}^{\infty} p_{ii} \pqty{k}.
    \end{align*}

    \begin{enumerate}
        \item If \(f_i = 1\), then \(V_i = \infty\) with probability \(1\), so \(i\) is recurrent, and \(\Expt_i \pqty{V_i} = \infty\), so \(\sum_{n \geq 0} p_{ii} \pqty{k} = \infty\).
        \item If \(f_i < 1\), then by \autoref{lem:return-probabilities-and-total-number-of-visits}, \(V_i\) is finite with probability \(1\), so \(V_i\) is infinite with probability \(0\), and hence \(i\) is transient, so \(\Expt_i \pqty{V_i} < \infty\), so \(\sum_{n \geq 0} p_{ii} \pqty{k} < \infty\).
    \end{enumerate}
\end{proof}

\begin{theorem}[Communicating States Share Recurrence]
    If \(x \communicatewith y\), then \(x\) is recurrent if and only if \(y\) is recurrent.
\end{theorem}

\begin{proof}
    Suppose \(x\) is recurrent. Since \(x \communicatewith y\), there are integers \(k, l\), such that \(p_{xy} \pqty{k} > 0\) and \(p_{yx} \pqty{l} > 0\). Then for any \(n \geq 0\),
    \begin{align*}
        \sum_{n \geq 0} p_{yy} \pqty{n} &\geq \sum_{n \geq 0} p_{yy} \pqty{l + n + k} \\
        &\geq \sum_{n \geq 0} p_{yx} \pqty{l} p_{xx} \pqty{n} p_{xy} \pqty{k} \\
        &= p_{yx} \pqty{l} p_{xy} \pqty{k} \sum_{n \geq 0} p_{xx} \pqty{n} \\
        &= \infty.
    \end{align*}
    so \(y\) is recurrent.
\end{proof}

\begin{corollary}[Recurrence is a Class Property]
    If \(A \subseteq I\) is a communicating class, then either all states in \(A\) are recurrent, or all states in \(A\) are transient.
\end{corollary}

\begin{theorem}[Recurrent Communicating Classes are Closed]
    If \(C \subseteq I\) is a recurrent communicating class, then it is closed.
\end{theorem}

\begin{proof}
    Suppose it is not closed, and we have some \(x \in C\) recurrent, with \(x \leadto y\) and \(y \notin C\). Then there is \(m\), such that \(p_{xy} \pqty{m} > 0\), so
    \[
        \Prob_x \pqty{V_x < \infty} \geq \Prob_x \pqty{X_m = y} > 0
    \]
    so \(x\) is transient, giving a contradiction.
\end{proof}